% \iffalse meta-comment
%
% Copyright (C) 2020-2021 by Antoine Missier <antoine.missier@ac-toulouse.fr>
% ----------------------------------------------------------------------
% 
% This file may be distributed and/or modified under the
% conditions of the LaTeX Project Public License, either version 1.3
% of this license or (at your option) any later version.
% The latest version of this license is in:
%
%    http://www.latex-project.org/lppl.txt
%
% and version 1.3 or later is part of all distributions of LaTeX 
% version 2005/12/01 or later.
%
% \fi
%
% \iffalse
%<*driver>
\ProvidesFile{schooldocs.dtx}
%</driver>
%<package>\NeedsTeXFormat{LaTeX2e}[2005/12/01]
%<package>\ProvidesPackage{schooldocs}
%<*package>
    [2021/09/18 v1.1 .dtx schooldocs file]
%</package>
%
%<*driver>
\documentclass{ltxdoc}
\usepackage[utf8]{inputenc}
\usepackage[T1]{fontenc}
\usepackage[english]{babel}
\usepackage{lmodern}
\usepackage{lipsum}
\usepackage{hyperref}
\usepackage{boxedminipage2e}
\usepackage{schooldocs}
%\usepackage{geometry}
\geometry{textwidth=355pt,vmarginratio=1:1,hmarginratio=3:2}
% modified by schoodocs loading, very close from the default geometry used by ltxdoc

\newcommand{\bs}{\textbackslash}
\setlength{\fboxsep}{13pt}
\newlength{\mpwd} \setlength{\mpwd}{5.25cm}
\newenvironment{stylex}{\begin{boxedminipage}[c][1.3\mpwd][t]{\mpwd}\tiny
    }{\end{boxedminipage}}
\newcommand*{\centered}[1]{{\setlength{\parskip}{0pt}\par\noindent\hfill
    #1\hfill\mbox{}}}
\providecommand\headtitlestyle{\scshape}
\providecommand\headsubjectstyle{\scshape}
\providecommand\schoolstyle{\scshape}

%\EnableCrossrefs         
%\CodelineIndex
%\RecordChanges

\title{The \textsf{schooldocs} package}%\thanks{This document
%    corresponds to \textsf{schooldocs}~\fileversion, dated \filedate.}}
\subject{Antoine Missier \\ % 
    \texttt{antoine.missier@ac-toulouse.fr}}
\date{\filedate, \fileversion}

\hypersetup{%
    colorlinks, 
    linkcolor=blue,
    pdftitle={schoodocs}, 
    pdfsubject={LaTeX package}, 
    pdfauthor={Antoine Missier}
    }

\begin{document}
  \DocInput{schooldocs.dtx}
  \PrintChanges
  \PrintIndex
\end{document}
%</driver>
% \fi
%
% \CheckSum{766}
%
% \CharacterTable
%  {Upper-case    \A\B\C\D\E\F\G\H\I\J\K\L\M\N\O\P\Q\R\S\T\U\V\W\X\Y\Z
%   Lower-case    \a\b\c\d\e\f\g\h\i\j\k\l\m\n\o\p\q\r\s\t\u\v\w\x\y\z
%   Digits        \0\1\2\3\4\5\6\7\8\9
%   Exclamation   \!     Double quote  \"     Hash (number) \#
%   Dollar        \$     Percent       \%     Ampersand     \&
%   Acute accent  \'     Left paren    \(     Right paren   \)
%   Asterisk      \*     Plus          \+     Comma         \,
%   Minus         \-     Point         \.     Solidus       \/
%   Colon         \:     Semicolon     \;     Less than     \<
%   Equals        \=     Greater than  \>     Question mark \?
%   Commercial at \@     Left bracket  \[     Backslash     \\
%   Right bracket \]     Circumflex    \^     Underscore    \_
%   Grave accent  \`     Left brace    \{     Vertical bar  \|
%   Right brace   \}     Tilde         \~}
%
%
% \changes{v0.0}{2010/09/01}{Initial personal version}
% \changes{v0.1}{2020/01/05}{Initial documentation in French}
% \changes{v0.2}{2020/05/04}{Examples file}
% \changes{v1.0}{2020/07/13}{Initial public version with documentation in English}
% \changes{v1.1}{2021/09/18}{Using the translations package for internationalization}
%
% \GetFileInfo{schooldocs.dtx}
% 
% \maketitle
% 
% \section{Introduction}
%
% The purpose of this package is to provide several layout styles for school documents.
% It can be used for exercises sheets, exams, course materials.
% The package sets the page geometry (dimensions of text and margins)
% and the title typesetting;
% the various \emph{styles} define the header, footer and title formatting, 
% Many features are freely configurable.
% Six different styles are proposed which are named:
% \begin{center} \bfseries
% \hfill classic \hfill elegant \hfill modern \hfill soft 
% \hfill exam \hfill collection \hfill \mbox{}
% \end{center}
% If the document is an exam subject to be completed and returned,
% we have a particular style, \textbf{identity}, in which
% left and center parts of the header are replaced by the words 
% \textsf{\identityname} and \textsf{\identityforename}.
%
% \medskip
% The styles are presented in paragraph \ref{par:styles}
% and the file \href{run:schooldocs-examples.pdf}{\textsf{schooldocs-examples.pdf}}
% shows the result of each.
%
% \section{Usage}
%
% \subsection{Choosing a style}
% After loading the package with |\usepackage{schooldocs}|,
% we define the style by placing |\pagestyle|\marg{style} in the preamble.
% This command comes from the package \textsf{fancyhdr} by Piet van Oostrum.
% It can be placed before of after entering information described in the next subsection.
%
% \medskip
% In the text body, you can change the style of a particular page by placing 
% |\thispagestyle|\marg{style} at the page you want.
% For example, you can place the command |\thispagestyle{identity}|
% in an appendix page to render.
% \subsection{Information fields}
%
% \DescribeMacro{\title}
% Usually a \LaTeX\ document title is made up of information
% provided by the macros
% |\title|, |\date| and |\author|.
% The |\title| macro has been redefined with an enriched syntax: 
% |\title|\oarg{head}\marg{title}.
% The optional parameter \meta{head} corresponds to the text displayed in the header
% instead of \meta{title}.
% When not provided, the mandatory argument \meta{title} will be used in headers.
%
% The macros |\date| and |\author| have not been redefined but the information provided
% is also used in titles, headers or footers
% (depending on the style used)\footnote{The macro \texttt{\bslash and},
% sometimes used in \texttt{\bslash author} 
% to typeset a title with several authors, doesn't work here, 
% except in the \textbf{collection} style. It produced a compilation error.
% The macro \texttt{\bslash thanks} also doesn't work but without producing error.}.
% 
% \pagebreak
% \medskip
% \DescribeMacro{\subject} \DescribeMacro{\school} \DescribeMacro{\institute}
% In addition to the previous ones, we provide two other information macros,
% used in the title, or the headers or the footers, depending on the style:
% |\subject| and |\school| (or its alias |\institute|).
% The macro |\subject| has an optional argument, like in |\title|: 
% |\subject|\oarg{head}\marg{title} which allows to place an alternative text
% in header, different from the one displayed in the title.
%
% \medskip 
% The information macros, that have to be called in the preamble\footnote{%
% Nevertheless, the macros \texttt{\bslash date} and \texttt{\bslash author}
% can be used in the document body.},
% are optional. If omitted the corresponding field will be empty,
% except for the date which will display today's date.
% So, if you do not want any date, just write |\date{}|.
% When title elements like subject, date are missing, the title height will be reduced accordingly.
%
% Except |\date|, all these information macros accept multi-line arguments,
% for example  |\school{Royal College of Pataphysics\\London}|.
% To typeset the main title on two (or more) lines, use |\\[1ex]|
% for a good vertical spacing.
%
% When using the \textsf{hyperref} package (by Sebastian Rahtz and Heiko Oberdiek)
% to produce a pdf output, information provided by |\title|, |\subject| and |\author| 
% (and |\school|\footnote{The \emph{author} field of the pdf properties will be filled
% like this: \textsf{author (school)},
% but \texttt{\bs author} should be called before \texttt{\bs school} in the preamble.
% })
% will be written in the pdf properties also.
%
% \medskip
% The file \textsf{schooldocs-examples.pdf} uses the examples of the following table.
% \begin{center}
% \begin{tabular}{ll}
% \hline
% macro & example \\
% \hline
% |\title| & |\title[The classic style]{The classic style layout}| \\
% |\subject| & |\subject[Pathography / M2A]{Pathography / Master 2A}| \\
% |\author| & |\author{Antoine Missier}| \\
% |\school| & |\school{Royal College of Pataphysics}| \\
% |\date| & |\date{May 4, 2020}| \\
% \hline
% \end{tabular}
% \end{center}
%
% \DescribeMacro{\subtitle}
% In the \textbf{exam} style, the title uses an additional field
% provided by the macro |\subtitle|\marg{text}. The parameter \meta{text}
% will be placed under the other title elements and before the separation rule,
% if there is one. For instance |\subtitle{Duration of the test: 2\,h}|.
% For a multi-line argument, the line break instruction must be
% |\par| and not |\\|.
%
% \subsection{Making the title}
%
% \DescribeMacro{\maketitle}
% As in \LaTeX\ standard classes, the title is produced 
% by the macro |\maketitle|, which should be placed after |\begin{document}|.
% This command has been redefined by the package.
% Indeed, the standard title produced by the \LaTeX\ class \texttt{article},
% seems inappropriate for school documents like an exercise sheet for example.
% The macro redefinition reduces the vertical space before the title
% and typesets an horizontal rule to separate the title from the following text.
% By default, the title is centered but we can change this, 
% as presented in section \ref{par:title-set}, to get 
% a left or right-aligned title (an also boxed).
% Various titles are presented in the file 
% \href{run:schooldocs-examples.pdf}{\textsf{schooldocs-examples.pdf}}.
% 
% The macro has on optional parameter |\maketitle|\oarg{rulelength}
% to set the length of the separation rule, by default 2.5\,cm.
% The rule can be totally removed by setting the
% optional argument to 0\,cm; we can also make a rule of
% whole line width with |\maketitle[\linewidth]| or |\maketitle[\textwidth]|.
%
% Each style has its own title formatting. Nevertheless the package provides
% a ``standard'' title when no style is loaded or with
% the \LaTeX\ styles \textbf{empty} or \textbf{plain}. 
% The \textbf{elegant} style keeps also this ``standard'' style.
%
% \medskip
% \DescribeMacro{\seprule}
% As in the title, a rule can also be used in the text body to make separations.
% This is done by the macro |\seprule|\oarg{length} which
% produces an horizontal centered rule of (optional) \meta{length}, by default 2.5\,cm.
% 
% \medskip
% \DescribeMacro{\correct}
% For typesetting a correct version of an exercise sheet, we provide the macro |\correct|,
% which adds \textsf{-- Correct Version} to the title. 
% It should be placed in the preamble.
% Combined with packages that deal with multiple versions of a document,
% this macro avoids to define different titles for a an exercise sheet an its correct version.
% However in the \textbf{exam} style, the text \textsf{Correct version}
% will not be added in the title (except in headers), but will be substituted for
% the subtitle.
% 
% \medskip
% \DescribeMacro{\makesmalltitle}
% As in the package \textsf{phfnote} by Philippe Faist, it seems wisely
% to propose a compact title, where information
% provided by |\title|, |\subject| and |\date| are presented
% on a single line (see \textsf{schooldocs-examples.pdf}).
% The rule is always displayed and expands on the whole line width.
% 
% \subsection{Style layouts} \label{par:styles}
%
% According to good typographic practice, the header on the title page should be empty
% (see The \LaTeX\ Companion).
% This ``rule'' is seldom observed in school documents,
% and, on the first page, we find often information (author, school, etc.)
% on top left and the date on top right, as for a letter.
% Except for the \textbf{classic} style, that we present first, 
% we wanted to comply with good practice,
% but we keep some information in the footer yet.
%
% We present below miniature views to show the layout of the different styles. 
% The file \href{run:schooldocs-examples.pdf}{\textsf{schooldocs-examples.pdf}}
% gives a real view of them.
% The rule (width) under the header or over the footer
% is freely configurable, as well as the format of each element and 
% also some space settings, see section \ref{par:head-settings}.
%
% \medskip
% \begin{stylex}
% {\schoolstyle School} \hfill {Author}\\[-0.3ex]
% {\headsubjectstyle Subject} \hfill {date}\\
% \vspace{0.04\mpwd}
% \centered{\footnotesize\sffamily\bfseries The classic style}
% \vspace{-0.25ex}\centered{\rule{0.15\mpwd}{0.1pt}}
% \vspace{0.85\mpwd}
% \centered{1/2}
% \end{stylex}
% \qquad
% \begin{stylex}
%  \mbox{}\\[1ex]
% {\headtitlestyle Title} \hfill {\headsubjectstyle Subject}\\
% \vspace{0.97\mpwd}
% \centered{2/2} 
% \end{stylex}
% \medskip
%
% \begin{stylex}
% \mbox{}\vspace{0.03\mpwd}
% \begin{center}
% {\footnotesize\sffamily\bfseries The elegant style}\\[1.7ex]
% {\scriptsize Subject} \\[0.5ex]
% {\scriptsize date} \\
% \vspace{-0.5ex}\rule{0.15\mpwd}{0.1pt}
% \end{center}
% \vspace{0.72\mpwd}
% {\schoolstyle School} \hfill Author
% \end{stylex}
% \qquad
% \begin{stylex}
% {\headtitlestyle Title} \hfill Page
% \end{stylex}
% \medskip
%
% \begin{stylex}
% \mbox{}\vspace{0.03\mpwd}
% \begin{center}
% {\footnotesize\sffamily\bfseries The modern style}\\[1.7ex]
% {\scriptsize Subject} \\
% \vspace{-0.5ex}\rule{0.15\mpwd}{0.1pt}
% \end{center}
% \vspace{0.75\mpwd}
% {\schoolstyle School} \hfill \begin{tabular}[c]{r@{}}Author\\[-0.5ex] date\end{tabular}
% \end{stylex}
% \qquad
% \begin{stylex}
% {\headtitlestyle Title} \hfill {\headsubjectstyle Subject} \par
% \vspace{-1.9ex}\rule{0.82\mpwd}{0.1pt}
%
% \vspace{1.01\mpwd}
% \centered{2/2}
% \end{stylex}
% \medskip
%
% \begin{stylex}
% \mbox{}\vspace{0.03\mpwd}
% \begin{center}
% {\footnotesize\sffamily\bfseries The soft style}\\
% \vspace{-0.5ex}\rule{0.15\mpwd}{0.1pt}
% \end{center}
% \vspace{0.83\mpwd}
% {\schoolstyle School} \hfill {\headsubjectstyle Subject}\\
% Author \hfill date
% \end{stylex}
% \qquad
% \begin{stylex}
% {\headtitlestyle Title} \hfill Page \par
% \vspace{-1.9ex}\rule{0.82\mpwd}{0.1pt}
% \vspace{0.97\mpwd}
%
% \schoolstyle School \hfill \headsubjectstyle Subject
% \end{stylex}
% \medskip
%
% \begin{stylex}
% \mbox{}\vspace{0.1\mpwd}
% \begin{center}
% {\footnotesize THE EXAM STYLE}\\[3.2ex]
% {\scriptsize\bfseries Subject} \\[2.8ex]
% {\scriptsize date} \\[2ex]
% {\scriptsize\slshape subtitle} \\
% \rule{0.15\mpwd}{0.1pt}
% \end{center}
% \vspace{0.51\mpwd}
% \begin{tabular}[c]{@{}l}Author\\[-0.5ex] \schoolstyle School\end{tabular} \hfill {Page}
% \end{stylex}
% \qquad
% \begin{stylex}
% {\headtitlestyle TITLE} \hfill {\headsubjectstyle Subject} \par
% \vspace{-1.9ex}\rule{0.82\mpwd}{0.1pt}
% 
% \vspace{1.01\mpwd}
% \begin{tabular}[c]{@{}l}Author\\[-0.5ex] \schoolstyle School\end{tabular} \hfill {Page}
% \end{stylex}
% \medskip
%
% \begin{stylex}
% \mbox{}\vspace{0.1\mpwd}
% \begin{center}
% {\footnotesize THE COLLECTION STYLE}\\[3.2ex]
% {\scriptsize\bfseries Subject} \\[2.8ex]
% {\scriptsize Author} \\[1.5ex]
% {\scriptsize School} \\[1.5ex]
% {\scriptsize date} \\
% \rule{0.15\mpwd}{0.1pt}
% \end{center}
% \end{stylex}
% \qquad
% \begin{stylex}
% {\headtitlestyle TITLE} \hfill {\headsubjectstyle Section} \par
% \vspace{-1.9ex}\rule{0.82\mpwd}{0.1pt}
% 
% \vspace{1.01\mpwd}
% {\schoolstyle School} \hfill {2}
% \end{stylex}
% \bigskip
%
% In the \textbf{collection} style, usefull for longer documents,
% the name of the current section appears in the right header 
% and changes as things progress
% as in the \LaTeX\ \textbf{book} class.
%
% \section{Settings}
%
% \subsection{Page layout}
%
% The page layout is set in the same way, for all styles.
% It is defined by the macro |\geometry| of the package
% \textsf{geometry} (by Hideo Umeki and David Carlisle).
% According to good typographic practice (see the \LaTeX\ Companion,
% the \textsf{typearea} package or \textsf{KOMA-Script} bundle documentation)
% the ratio between the text body (the type area as named by typographers)
% and the page size must be the same horizontally and vertically. 
% This ratio gauges the ``shades of gray'' of a page.
% The default scale in the \textsf{geometry} package is 0.7.
% For school documents, it seems that this setting leaves too much white space in the margins, thus this ratio has been set to 0.75.
% Of course, it can be freely redefined with the |\geometry| command.
% Furthermore, according to ``canons of page construction'', in a one-sided document,
% the ratio between the left and right margins should be 1:1,
% while the ratio between the top and bottom margins should be 1:2.
% It seems that this recommended ratio of 1:2 is seldom observed
% and it produces a very wide bottom margin.
% Besides, when loading the \textsf{geometry} package, the default vertical ratio is 2:3 and not 1:2.
% To get more balanced margins, we have prefer to set it at 3:4:
% |\geometry{scale=0.75,vmarginratio=3:4,heightrounded}|\footnote{The option
% \texttt{heightrounded} rounds text height to $n$-times ($n$ an integer) the height of a line, to avoid ``underfull vbox'' in some cases.}.
%
% \subsection{Title settings} \label{par:title-set}
%
% Several parameters allow to customize the title composition.
% They are presented in the following table with their default value.
% Colors are changed with |\definecolor|, macros with |\renewcommand|
% and lengths with |\setlength|.
% Most of command names are explicit enough to give some information about their use.
% 
% Let us specify nevertheless that |\titleflush| manages the alignment: 
% centered by default, it can be left-aligned with
% |\renewcommand{\titleflush}{flushleft}| or right-aligned with
% |\renewcommand{\titleflush}{flushright}|\footnote{This macro has no effect 
% with the \textbf{exam} and \textbf{collection} styles, 
% the title is always centered.}.
%
% The length |\titlesep| sets the vertical spacing between the fields \emph{title}
% and \emph{subject}, except for the \textbf{classic} and \textbf{soft} styles,
% witch contain only the field \emph{title}, or if the field \emph{subject}
% is not given. 
% Vertical spaces between other title fields are defined in relation to
% |\titlesep| (|0.5\titlesep| for most).
%
% Unless otherwise stated, the column ``special settings'' of the following table
% applies to the \textbf{exam} and \textbf{collection} styles.
%
% \begin{center} \small
% \begin{tabular}{ccc}
% \hline
% parameter & default settings & special settings \\
% \hline
% |titlecolor| & |black| & \\
% |\titlestyle| & |\LARGE\sffamily\bfseries| & |\LARGE\MakeUppercase| \\
% |\subjectstyle| & |\large| & |\Large\bfseries|\\
% |\datestyle| &  |\large| \\
% |\smalltitledatestyle| & |\large\slshape| & \\
% |\titleflush| & |center| & \\
% |\titletopskip| & -1.32\,cm  & -0.67\,cm (\textbf{classic}), 0.67\,cm \\
% |\titlebottomskip| & |\medskipamount| \\
% |\titlesep| & |2\medskipamount| & |2\bigskipamount| \\
% |\seprulewidth| & 0.3\,pt &\\
% |\seprulelength| & 2.5\,cm  &\\
% \hline
% \end{tabular}
% \end{center}
%
% \DescribeMacro{\subtitlestyle}\DescribeMacro{\titlecorrectstyle}
% The \textbf{exam} style has two additional setting macros: 
% |\subtitlestyle|, set by default to  |\large\slshape|,
% and |\titlecorrectstyle|, for the composition of \textsf{Correct Version}
% produced by |\correct|, set by default to |\color{red}\large\scshape|.
%
% In the \textbf{collection} style, the fields
% \emph{school} and \emph{author} are included into the title. Their features are then 
% controlled by |\authorstyle|\footnote{For other styles, the macro
% \texttt{\bs authorstyle} belongs to headers or footers parameters.} set to |\Large|.
%
% If we want to create a separate page for the title, in particular for the 
% \textbf{exam} or \textbf{collection} styles, we can place |\maketitle| into
% an environment |titlepage| provided by \LaTeX.
%
% \medskip
% \DescribeMacro{\boxedshape}
% A particular formatting has been predefined to produce a slightly shaded box around the title: |\boxedshape|\marg{text}
% (thanks to the \textsf{fancybox} package by Timothy Van Zandt).
% It is not activated by default but it can be used within the
% command |\titlestyle| if we want to get this kind of effect, e.g.
% |\renewcommand{\titlestyle}{\LARGE\sffamily\bfseries\boxedshape}|.
%
% \subsection{Headers and footers} \label{par:head-settings}
%
% Headers and footers typeset is produced thanks to 
% dedicated macros\footnote{The \textsf{fancyhdr} package
% provides the commands \texttt{\bs lhead}, \texttt{\bs chead}, \texttt{\bs rhead}
% for the left, center and right parts of the header, and the same for the footer:
% \texttt{\bs lfoot}, \texttt{\bs cfoot} and \texttt{\bs rfoot}.}
% of the \textsf{fancyhdr} package.
% As for the title, several macros allow to customize certain settings.
% They are presented in the table below.
% Let us specify that |\pagenamestyle| defines the format of the word \emph{page} 
% used in front of the page number.
%
% The |\headstyle| and |\footstyle| commands are generic for all elements 
% of headers and footers. They are combined with macros giving
% specific settings for each field.
% \begin{center}
% \begin{tabular}{cc}
% \hline
% parameter & default setting \\
% \hline
% |headingcolor| & |black| \\
% |\headstyle| & |\small\color{headingcolor}| \\
% |\footstyle| & |\small\color{headingcolor}| \\
% |\headtitlestyle| & |\scshape| \\
% |\headsubjectstyle| & |\scshape| \\
% |\schoolstyle| & |\scshape| \\
% |\headdatestyle| & |\textnormal| \\
% |\authorstyle| & |\textnormal| \\
% |\pagenamestyle| & |\textnormal| \\
% \hline
% \end{tabular}
% \end{center}
%
% The \textbf{exam} and \textbf{collection} styles have a specific setting for
% |\headtitlestyle|, defined by |\renewcommand{\headtitlestyle}{\MakeUppercase}|.
%
% \medskip
% Macros that allow to display (or remove) the rule in the header\footnote{Indeed 
% we have redefined the \texttt{\bs headrule} macro
% such that the rule never appears
% on the first page.}
% or the footer
% come from the \textsf{fancyhdr} package. For example, in any style,
% we remove the rule in the header with:
% |\renewcommand{\headrulewidth}{0pt}|,
% and display it in the footer with:
% |\renewcommand\footrulewidth{0.4pt}|
%
% \subsection{Internationalization}
% Four keywords are used by the package: \emph{page}, \emph{Correct version}, 
% \emph{Name} and \emph{Forename},
% defined by the macros |\pagename|, |\correctname|, |\identityname| and |\identityforename|.
% |\pagename| is used in header or footer, |\correctname| by the macro |\correct|
% and both others by the \textbf{identity} style.
% These keywords get automatic translation in a few languages\footnote{Translation
% is currently integrated into the package for the following languages:
% French, German, Spanish, Italian, Portuges.}
% according to the selected language by \textsf{babel} package or \textsf{polyglossia}.
% They can be redefined, with |\renewcommand|,
% which should be placed \emph{after} |\begin{document}|,
% or else we can use macros from the \textsf{translations} package by Clemens Niederberger
% (allowing dynamic language change), e.g.\\
% |\DeclareTranslation{romanian}{pagin\u{a}}| (in the preamble).
%
%
% \StopEventually{}
%
% \section{Implementation}
%
% \subsection{Required packages and geometry settings}
%
%    \begin{macrocode}
\ProvidesPackage{schooldocs}
\NeedsTeXFormat{LaTeX2e}

\RequirePackage{geometry}
\let\footruleskip\relax % because memoir class uses it
\RequirePackage{fancyhdr}
\RequirePackage{ifthen}
\RequirePackage{lastpage}
\RequirePackage{fancybox}
\RequirePackage{xcolor}
\RequirePackage{translations}

\geometry{scale=0.75,vmarginratio=3:4,heightrounded}

%    \end{macrocode}
% It is recommended to define the settings with the command |\geometry|
% rather than as option when loading the package \textsf{geometry}, to avoid option clash.
% 
% To change the scale factor, if the |\geometry| command doesn't work, use |\newgeometry|.
%\medskip
% \subsection{Internationalization}
%
%    \begin{macrocode}
\DeclareTranslationFallback{schooldocs-pagename}{page}
\DeclareTranslationFallback{schooldocs-correct}{Correct version}
\DeclareTranslationFallback{schooldocs-name}{Name:}
\DeclareTranslationFallback{schooldocs-forename}{Forename:}

\DeclareTranslation{english}{schooldocs-pagename}{page}
\DeclareTranslation{english}{schooldocs-correct}{Correct version}
\DeclareTranslation{english}{schooldocs-name}{Name:}
\DeclareTranslation{english}{schooldocs-forename}{Forename:}

\DeclareTranslation{french}{schooldocs-pagename}{page}
\DeclareTranslation{french}{schooldocs-correct}{Corrigé}
\DeclareTranslation{french}{schooldocs-name}{Nom :}
\DeclareTranslation{french}{schooldocs-forename}{Prénom :}

\DeclareTranslation{german}{schooldocs-pagename}{Seite}
\DeclareTranslation{german}{schooldocs-correct}{Musterlösung}
\DeclareTranslation{german}{schooldocs-name}{Name:}
\DeclareTranslation{german}{schooldocs-forename}{Vorname:}

\DeclareTranslation{spanish}{schooldocs-pagename}{página}
\DeclareTranslation{spanish}{schooldocs-correct}{Correccion}
\DeclareTranslation{spanish}{schooldocs-name}{Apellido:}
\DeclareTranslation{spanish}{schooldocs-forename}{Nombre:}

\DeclareTranslation{italian}{schooldocs-pagename}{pagina}
\DeclareTranslation{italian}{schooldocs-correct}{Corretto}
\DeclareTranslation{italian}{schooldocs-name}{Nome:}
\DeclareTranslation{italian}{schooldocs-forename}{Conome:}

\DeclareTranslation{portuges}{schooldocs-pagename}{página}
\DeclareTranslation{portuges}{schooldocs-correct}{Corrigiu}
\DeclareTranslation{portuges}{schooldocs-name}{Nome:}
\DeclareTranslation{portuges}{schooldocs-forename}{Primeiro nome:}

\newcommand*\pagename{\GetTranslation{schooldocs-pagename}}
\newcommand*\correctname{\GetTranslation{schooldocs-correct}}
\newcommand*\identityname{\GetTranslation{schooldocs-name}}
\newcommand*\identityforename{\GetTranslation{schooldocs-forename}}

%    \end{macrocode}
%
% \subsection{Information fields}
%
% After being defined, information fields are set to empty,
% to make tests with them work well.
%
%    \begin{macrocode}
\renewcommand{\title}[2][]{\long\def\@title{#2} 
    \ifthenelse{\equal{#1}{}}{\def\@headtitle{#2}}{\def\@headtitle{#1}}
    \@ifpackageloaded{hyperref}{\hypersetup{pdftitle={\@headtitle}}}{} 
} 
\title{}
\author{}

\newcommand\school[1]{\long\def\@school{#1}
    \@ifpackageloaded{hyperref}{
        \hypersetup{pdfauthor={\@author\ (\@school)}}}{}
} 
\school{}
\newcommand{\institute}{\school}

\let\subject\relax
\newcommand{\subject}[2][]{\long\def\@subject{#2} 
    \ifthenelse{\equal{#1}{}}{\def\@headsubject{#2}}{\def\@headsubject{#1}}
    \@ifpackageloaded{hyperref}{\hypersetup{pdfsubject={\@headsubject}}}{}
}
\def\@subject{}
\def\@headsubject{}
%    \end{macrocode}
% Because \textsf{KOMA-Script} package also uses the |\subject| macro,
% we make a |\relax| on this command.
%
% |\subject{}| doesn't work with the |\@empty| test used below in |\maketitle|, thus
% we define |\@subject| and |\@headsubject| as empty by default.
%
% \subsection{Settings}
%
% First we present the title settings.
%    \begin{macrocode}
\definecolor{titlecolor}{named}{black} 
\newcommand{\titlestyle}{\LARGE\sffamily\bfseries}
\newcommand{\subjectstyle}{\large}
\newcommand{\datestyle}{\large}
\newcommand{\smalltitledatestyle}{\large\slshape}
\newcommand{\titleflush}{center}
\newlength{\titletopskip} \setlength{\titletopskip}{-1.32cm}
\newlength{\titlesep} \setlength{\titlesep}{2\medskipamount}
\newlength{\titlebottomskip} \setlength{\titlebottomskip}{\medskipamount}
\newlength{\seprulewidth} \setlength{\seprulewidth}{0.3pt}
\newlength{\seprulelength} \setlength{\seprulelength}{2.5cm}

\newcommand{\boxedshape}[1]{%
    \color{gray} % no effect
    \setlength{\fboxsep}{4.5pt} 
    \setlength{\shadowsize}{3pt}
    \shadowbox{\quad\textcolor{titlecolor}{\rule{0pt}{2ex}#1}\quad} 
}

%    \end{macrocode}
% Default value for |\fboxsep| is 3\,pt and for |\shadowsize| is 4\,pt.
% These settings are local.
%
% \medskip
% Next the headers and footers settings.
%    \begin{macrocode}
\definecolor{headingcolor}{named}{black}
\newcommand{\headstyle}{\small\color{headingcolor}}
\newcommand{\footstyle}{\small\color{headingcolor}}
\newcommand{\headtitlestyle}{\scshape}
\newcommand{\headsubjectstyle}{\scshape}
\newcommand{\schoolstyle}{\scshape}
\newcommand{\headdatestyle}{\textnormal}
\newcommand{\authorstyle}{\textnormal}
\newcommand{\pagenamestyle}{\textnormal}
\setlength{\headheight}{\baselineskip} 
\renewcommand{\headrule}{\ifthenelse{\thepage=1}%
    {\vspace{-2ex}\rule{\headwidth}{0pt}}%
    {\vspace{-2ex}\rule{\headwidth}{\headrulewidth}}
}
    
%    \end{macrocode}
% The |\headrule| macro redefinition allow to remove the rule on the first page.
%
% The header height is set to |\baselineskip| to avoid
% error messages (``headheight too small'') when the text size is 
% greater than or equal to |\small|.
%
% \subsection{Title composition}
%
%    \begin{macrocode}
\renewcommand{\maketitle}[1][\seprulelength]{
    \mbox{}\par\vspace{\titletopskip}
    \begin{\titleflush}
        {\color{titlecolor}
        {\titlestyle{\@title}}
        \ifx\@subject\@empty\else
            \par\vspace{\titlesep}\subjectstyle\@subject
        \fi
        \ifthenelse{\equal{\@date}{}}{}{%
            \par\vspace{0.5\titlesep} \datestyle\@date}
        \par\vspace{0.5\titlesep}
        }
        \ifthenelse{\lengthtest{#1 > 0cm}}{\rule{#1}{\seprulewidth}}{}
    \end{\titleflush}
    \vspace{\titlebottomskip} 
}

\newcommand{\makesmalltitle}{
    \mbox{}\par\vspace{1.5\titletopskip}
    \begin{center}
        {\color{titlecolor}
        {\titlestyle{\@title}}
        \hfill \subjectstyle\@subject
        \ifthenelse{\equal{\@date}{}}{}{%
            \enskip --\enskip \smalltitledatestyle\@date}
        }
        \rule{\linewidth}{\seprulewidth}
    \end{center}
    \vspace{\titlebottomskip}
}

\newcommand\correct{%
    \let\@originalheadtitle\@headtitle
    \def\@headtitle{\@originalheadtitle ~--~\correctname}
    \@ifpackageloaded{hyperref}{%
        \AtBeginDocument{\hypersetup{pdftitle={\@headtitle}}}}{}
    \let\@originaltitle\@title
    \def\@title{\@originaltitle ~--~\correctname}
}

\newcommand{\seprule}[1][\seprulelength]{\begin{center}
    \raisebox{0.25\baselineskip}{\rule{#1}{\seprulewidth}}\end{center}}
    
%    \end{macrocode} 
% |\AtBeginDocument| is necessary in the |\correct| macro because 
%|\correctname| can be redefined by \textsf{babel} at |\begin{document}|.
%
% \medskip
% \subsection{The basic styles}
%
% The header in the \textbf{classic} style is displayed on two lines, therefore we have to increase the height
% |\headheight|.
%
% If we want to include header or footer in the total body,
% we can add the command |\geometry{includehead}| respectively 
% |\geometry{includefoot}| in the preamble.
%    \begin{macrocode}
\fancypagestyle{classic}{%
    \renewcommand{\headrulewidth}{0pt}
    \addtolength{\headheight}{\baselineskip}
    \lhead{\headstyle\ifthenelse{\thepage=1}{\schoolstyle\@school\mbox{}\\
        \headsubjectstyle\@subject}{\headtitlestyle\@headtitle} } 
    \rhead{\headstyle\ifthenelse{\thepage=1}{\authorstyle\@author\mbox{}\\ 
        \headdatestyle\@date}{\headsubjectstyle\@headsubject}}
    \cfoot{\footstyle \thepage /\pageref{LastPage}}
    \setlength{\titletopskip}{-0.67cm}
    \renewcommand{\maketitle}[1][\seprulelength]{%
        \mbox{}\par\vspace{\titletopskip}
        \begin{\titleflush}
            {\color{titlecolor}\titlestyle{\@title}}
            \par\vspace{0.5\titlesep}
            \ifthenelse{\lengthtest{##1>0cm}}{\rule{##1}{\seprulewidth}}{}            
        \end{\titleflush}
        \vspace{\titlebottomskip} 
    }    
}

\fancypagestyle{elegant}{%
    \renewcommand{\headrulewidth}{0pt}
    \lhead{\headstyle
        \ifthenelse{\thepage=1}{}{\headtitlestyle\@headtitle}}
    \rhead{\headstyle
        \ifthenelse{\thepage=1}{}{%
        \pagenamestyle\pagename~\thepage /\pageref{LastPage}}}
    \lfoot{\ifthenelse{\thepage=1}{\footstyle\schoolstyle\@school}{}}
    \cfoot{} 
    \rfoot{\ifthenelse{\thepage=1}{\footstyle\authorstyle\@author}{}}
}

\fancypagestyle{modern}{%
    \lhead{\headstyle
        \ifthenelse{\thepage=1}{}{\headtitlestyle\@headtitle}}
    \rhead{\headstyle
        \ifthenelse{\thepage=1}{}{\headsubjectstyle\@headsubject}}
    \lfoot{\ifthenelse{\thepage=1}{\footstyle\schoolstyle\@school}{}}
    \cfoot{\ifthenelse{\thepage=1}{}{
        \footstyle\thepage /\pageref{LastPage}}}
    \rfoot{\footstyle\ifthenelse{\thepage=1}{%
        \authorstyle\@author
        \ifx\@author\empty\else\\ \fi
        \headdatestyle\@date\ignorespaces}{}
    }
    \renewcommand{\maketitle}[1][\seprulelength]{% 
        \mbox{}\par\vspace{\titletopskip}
        \begin{\titleflush}
            {\color{titlecolor}
            {\titlestyle{\@title}}
            \ifx\@subject\@empty\else
                \par\vspace{\titlesep}\subjectstyle\@subject
            \fi
            \par\vspace{0.5\titlesep}
            }
            \ifthenelse{\lengthtest{##1>0cm}}{\rule{##1}{\seprulewidth}}{}
        \end{\titleflush}
        \vspace{\titlebottomskip} 
    }
}

\fancypagestyle{soft}{%
    \lhead{\headstyle
        \ifthenelse{\thepage=1}{}{\headtitlestyle\@headtitle}}
    \rhead{\headstyle
        \ifthenelse{\thepage=1}{}{%
        \pagenamestyle\pagename~\thepage /\pageref{LastPage}}}
    \lfoot{\footstyle{\schoolstyle\@school}
        \ifthenelse{\thepage=1}{\ifx\@author\empty\else\\ \fi
            \authorstyle\@author}{}
    }
    \cfoot{}
    \rfoot{\footstyle{\headsubjectstyle\@headsubject}
        \ifthenelse{\thepage=1}{\ifx\@date\empty\else\\ \fi
            \headdatestyle\@date\ignorespaces}{}
    }
    \setlength{\titletopskip}{-0.67cm}
    \renewcommand{\maketitle}[1][\seprulelength]{%
        \mbox{}\par\vspace{\titletopskip}
        \begin{\titleflush}
            {\color{titlecolor}\titlestyle{\@title}}
            \par\vspace{0.5\titlesep}
            \ifthenelse{\lengthtest{##1>0cm}}{\rule{##1}{\seprulewidth}}{}            
        \end{\titleflush}
        \vspace{\titlebottomskip} 
    }    
}

\fancypagestyle{identity}{%
    \lhead{\headstyle\headtitlestyle \identityname} 
    \chead{\headstyle\headtitlestyle \identityforename}
    \rhead{}
    }
    
%    \end{macrocode}       
%
% \subsection{The extended styles}
%
%    \begin{macrocode}
\fancypagestyle{exam}{%
    \let\subtitle\relax % because KOMA-Script uses it
    \newcommand\subtitle[1]{\long\def\@subtitle{##1}}
    \subtitle{}
    \renewcommand{\headtitlestyle}{\MakeUppercase}
    \lhead{\headstyle
        \ifthenelse{\thepage=1}{}{\headtitlestyle{\@headtitle}}}
    \rhead{\headstyle
        \ifthenelse{\thepage=1}{}{\headsubjectstyle\@headsubject}}
    \lfoot{\footstyle
        \authorstyle\@author
        \ifx\@author\empty\else\\ \fi
        \schoolstyle\@school
    }
    \cfoot{}        
    \rfoot{\footstyle \pagenamestyle\pagename~\thepage/\pageref{LastPage}}
    \setlength{\titletopskip}{0.67cm}
    \setlength{\titlesep}{2\bigskipamount}
    \renewcommand{\titlestyle}{\LARGE\MakeUppercase}
    \renewcommand{\subjectstyle}{\Large\bfseries}
    \newcommand{\subtitlestyle}{\large\slshape}
    \newcommand*{\titlecorrectstyle}{\color{red}\large\scshape}
    \renewcommand\correct{%
        \let\@originalheadtitle\@headtitle
        \def\@headtitle{\@originalheadtitle ~--~\correctname}
        \@ifpackageloaded{hyperref}{%
            \AtBeginDocument{\hypersetup{pdftitle={\@headtitle}}}}{}
        \renewcommand{\subtitlestyle}{\titlecorrectstyle}
        \subtitle{\correctname}
    }
    \renewcommand{\maketitle}[1][\seprulelength]{%
        \mbox{}\par\vspace{\titletopskip}
        \begin{center}
            {\color{titlecolor}
            {\titlestyle{\@title}}
            \par\vspace{\titlesep}{\subjectstyle \@subject}
            \par\vspace{\titlesep}
            \ifthenelse{\equal{\@date}{}}{}{%
                \datestyle\@date\par\vspace{0.5\titlesep}}
            \ifthenelse{\equal{\@subtitle}{}}{}{%
                \subtitlestyle\@subtitle\par\vspace{0.5\titlesep}
            }
            }
            \ifthenelse{\lengthtest{##1>0cm}}{\rule{##1}{\seprulewidth}}{}                 
        \end{center}
        \vspace{\titlebottomskip}
    }    
}

\fancypagestyle{collection}{%
    \renewcommand{\headtitlestyle}{\MakeUppercase}
    \lhead{\headstyle\headtitlestyle\@headtitle}
    \AtBeginDocument{\renewcommand{\sectionmark}[1]{\markboth{##1}{}}}
    \rhead{\headstyle\headsubjectstyle\leftmark} %\@section
    \lfoot{\footstyle\schoolstyle \@school}
    \cfoot{}
    \rfoot{\footstyle\thepage}
    \setlength{\titletopskip}{0.67cm}
    \setlength{\titlesep}{2\bigskipamount}
    \renewcommand{\titlestyle}{\LARGE\MakeUppercase}
    \renewcommand{\subjectstyle}{\Large\bfseries}
    \renewcommand{\authorstyle}{\Large}
    \renewcommand{\datestyle}{\Large}
    \renewcommand{\maketitle}[1][\seprulelength]{%
        \thispagestyle{empty}
        \mbox{}\par\vspace{\titletopskip}
        \begin{center}
            {\color{titlecolor}
            {\titlestyle{\@title}}
            \par\vspace{\titlesep} {\subjectstyle{\@subject}}
            \par\vspace{\titlesep} 
            \ifx\@author\empty
            \else \authorstyle
                \begin{tabular}[t]{c}\@author\end{tabular}
                \ifx\@school\empty\else\par\vspace{1ex} \@school\fi
                \par\vspace{0.5\titlesep}
            \fi
            \ifthenelse{\equal{\@date}{}}{}{%
                \datestyle\@date \par\vspace{0.5\titlesep}}
            }
            \ifthenelse{\lengthtest{##1>0cm}}{\rule{##1}{\seprulewidth}}{}     
        \end{center}
        \vspace{\titlebottomskip}     
    }
}
%    \end{macrocode}
%
% In the \textbf{collection} style, 
% we redefine |\sectionmark| with |\AtBeginDocument| to prevent |\MakeUppercase| 
% from eliminating the section number 
% (or you can use |\nouppercase| here).
%
% |\ifthenelse| doesn't work when |\and| is used within |\author|.
%
% Notice the |\pagestyle{empty}| in |\maketitle| which prevents
% any header on the first page.
%
% \Finale
\endinput
